% STATUS: OUTLINE ONLY. The rubric awards zero points for this section at
% the midterm, so it exists to show the structure. Every number below is
% already measured and recorded; what is missing is the prose and the
% figures, which is Thursday and Friday's work.
%
% Every figure listed here has a run directory behind it. Take the number
% from the run directory, never from this file, and never from memory.

\section{Evaluation}

% ---------------------------------------------------------------------
\subsection{Questions}
% Four, in order:
%   Q1  Does placement matter at all on a real cluster?
%   Q2  How much does following the traffic beat doing nothing?
%   Q3  How does that advantage change as the workload moves faster?
%   Q4  Does a learned policy beat the heuristic?
% Q4's answer is "not yet", and the section is structured so that is
% reported rather than buried.

% ---------------------------------------------------------------------
\subsection{Setup}
% Six regions, replication factor 3, 256 cells, 200 decision windows per
% episode, eight seeds. Round-trip matrix stated as a table.
% Cluster: six CockroachDB nodes in containers, netem latency injection
% accurate to 0.1 ms across the measured pairs.
% Hardware: GitHub Actions ubuntu-latest, 2 cores, 7 GB. State this.

% ---------------------------------------------------------------------
\subsection{Q1: placement changes latency on a real cluster}
% Source: Cluster pilot #2, 20 October 2026, docs/results-log.md run 2.
%
%                        p50      p95      stdev
%   leases co-located    272.7    275.5    48.8
%   leases spread       1013.4   1648.0   302.0
%   p95 gap 1372.5 ms
%
% THREE THINGS THAT MUST APPEAR IN THIS SUBSECTION:
%
% 1. "Co-located" means the leases, not the data. With replication factor
%    3 and one node per region, three replicas cannot fit in one region,
%    so every write pays a cross-region quorum even in the co-located
%    case. That is why the co-located p50 is 272.7 ms and not single
%    digits. The experiment isolates the leaseholder effect, which is
%    exactly what MOVE_LEASE corresponds to.
% 2. The confidence interval is indicative, not a precision claim. At
%    N = 200 transactions the p95 is the 190th order statistic, so the
%    bootstrap interval is degenerate at the top. Report the gap, report
%    N, say the interval is indicative, and state that choosing N by
%    measurement is outstanding work.
% 3. Six database processes shared two cores. Say so.

% ---------------------------------------------------------------------
\subsection{Q2: following the traffic beats doing nothing}
% Source: Experiments #1, 20 October 2026, 8 seeds, volatility 1.0.
%
%   Static hash                      156.11  [148.92, 164.33]     0.0 GB
%   Random                           158.02  [150.98, 165.63]   204.8 GB
%   Follow traffic                    98.66  [ 89.49, 106.58]   785.2 GB
%   Best per window (fixed shards)    93.66  [ 86.60,  99.63]     0.0 GB
%
% 36.8% improvement. Note that random is no better than static despite
% moving 205 GB, which is the point of including it: the gain comes from
% choosing where, not from moving.
%
% TABLE 1 goes here.

% ---------------------------------------------------------------------
\subsection{Q3: the advantage collapses as the workload moves faster}
% Source: Sweep #1, 20 October 2026, 8 seeds per value.
% THIS IS THE SECTION THE PAPER IS FOR. Give it the space.
%
%   volatility   static   follow   bound    headroom   advantage
%        0.5     156.27    94.31   93.87     0.44 ms      39.7%
%        1.0     156.11    98.66   93.66     5.00 ms      36.8%
%        2.0     155.37   125.37   89.47    35.90 ms      19.3%
%        4.0     155.77   139.58   82.38    57.20 ms      10.4%
%
%   GB moved by follow traffic: 405 -> 785 -> 1057 -> 1115
%
% The three observations, in this order:
%
%   (a) Static hashing is flat. It does not degrade with volatility
%       because it was never tracking anything. Say this plainly; it is
%       the reason the advantage column falls rather than the baseline
%       rising.
%   (b) The bound IMPROVES with volatility, 93.87 -> 82.38, because
%       sharper traffic concentration means a well-chosen placement does
%       better. So the achievable optimum gets better whilst the
%       heuristic gets worse. That widening gap is the finding.
%   (c) The heuristic pays more and gets less: movement rises from
%       405 GB to 1115 GB whilst its advantage falls from 39.7% to
%       10.4%. A reactive rule under fast drift is buying churn.
%
% The conclusion to draw: the room for a learned policy is largest
% exactly where the problem is hardest. At volatility 0.5 there is
% nothing to win. At volatility 4 there is 41% of the heuristic's own
% latency to win.
%
% FIGURE 2 goes here, two panels: latency against volatility, and GB
% moved against volatility. Source: results/sweeps/volatility/sweep.png.

% ---------------------------------------------------------------------
\subsection{Q4: the learned policy does not yet match the heuristic}
% Sources: Train agent #1 and #2, 20 October 2026, both at volatility 1.0,
% both evaluated on held-out seeds 2000-2007.
%
%   run 1, w_migration = 0.30:   146.55 ms,  53.2 GB moved
%   run 2, w_migration = 0.00:   135.57 ms, 192.8 GB moved
%   follow traffic:               90.76 ms, 752.6 GB moved
%
% Report this straight. The structure of the argument:
%
%   (a) Run 1 learned to sit still: 53 GB against the heuristic's 753.
%       The migration penalty is immediate and certain whilst the latency
%       benefit is delayed and spread over later windows, and a Q-function
%       trained for 1,500 episodes resolves that disagreement by not
%       moving.
%   (b) Run 2 removed the penalty entirely. Movement rose to 193 GB and
%       latency improved by 11 ms, which confirms (a) as a contributing
%       cause but does not close the gap.
%   (c) Run 2's held-out evaluation was still improving at the final
%       training step, 127.66 ms at step 300,000 against run 1's plateau
%       at 138 ms by step 160,000. The budget ran out before convergence.
%   (d) Both runs trained at volatility 1.0, where the sweep shows the
%       available headroom is 5 ms. The agent was asked to win a contest
%       with almost no prize. Training at volatility 2 to 4 is the
%       obvious next experiment and is outstanding work.
%
% FIGURE 3: training curve, held-out latency against timestep, both runs.
% Source: eval_history.json in each model directory.

% ---------------------------------------------------------------------
\subsection{Threats to validity}
% - The simulator's cost model is analytic. Q1 validates that placement
%   matters on real hardware; it does not validate the magnitudes.
% - Two cores for six database processes adds scheduling noise.
% - N = 200 transactions was chosen by convenience rather than by
%   measurement.
% - The bound holds the starting shard boundaries, so it is not a global
%   optimum and should not be read as one.
% - The agent selects inapplicable actions in 80-90% of windows. Cite
%   \cite{huang2022masking} for why masking matters as the invalid share
%   grows, and \cite{zahavy2018elimination} for an action-elimination
%   network built on DQN specifically, which is the closest published fix.
% - Results at volatility 1.0 should not be generalized to other
%   volatilities, which is precisely what Q3 demonstrates.
